\section{Experimental Plan}

This section is a forward-looking placeholder. Before the method can support
publication-quality inferential claims, the manuscript should cover at least
the following empirical blocks.

\subsection{Null calibration}

The first requirement is a null calibration study on binary data with no true
cluster structure. The goal is to measure how often Gate 2 and Gate 3 reject
under scenarios where they should not. This should be reported both at the
individual-test level and at the resulting number-of-clusters level. The latter
is critical because a method can have mild per-test inflation yet still produce
severe over-splitting once the tree walk is applied.

\subsection{Power under planted structure}

The second requirement is a controlled power study with planted binary subtrees.
This should vary:

\begin{itemize}[leftmargin=1.5em]
  \item sample size per cluster,
  \item effect size in feature rates,
  \item number of informative features,
  \item tree imbalance, and
  \item background noise level.
\end{itemize}

Power should be measured both for edge recovery and for final cluster recovery.
Adjusted Rand Index and exact split recovery are both useful endpoints.

\subsection{Ablations}

The current codebase makes several choices that can be ablated directly:

\begin{itemize}[leftmargin=1.5em]
  \item sibling calibration method,
  \item projection-dimension rule,
  \item edge multiple-testing correction,
  \item passthrough traversal on versus off, and
  \item optional post-hoc merge.
\end{itemize}

This section should eventually explain not only which configuration performs
best, but also which choices are required for stable Type I error and which are
mostly about power or cluster granularity.

\subsection{Real-data benchmarks}

The repository already contains benchmark infrastructure and reporting code. A
manuscript-quality benchmark section should align those experiments with the
method claims in the text, focusing on:

\begin{itemize}[leftmargin=1.5em]
  \item synthetic binary benchmarks with known labels,
  \item sensitivity to over-splitting and under-splitting,
  \item comparison to standard clustering baselines, and
  \item qualitative inspection of annotated trees.
\end{itemize}
